\section{Introduction}
\label{sec:intro}

\subsection{Problem Description}

The Traveling Salesman Problem (TSP) represents one of the most studied problems in combinatorial optimization, appealing greatly to modern 
logistics. The problem statement is deceptively simple: given a list of cities and the distances between any pair of cities, what is the 
shortest possible route that visits each city exactly once and returns to the origin city? Several algorithms have come up across decades, some 
that guarantee an optimal solution for smaller problems, while others produce near-optimal answers fast enough for real-life applications.

Formally, the problem is defined on a weighted complete graph $G = (V, E, w)$, where $V$ is the set of nodes (cities), $E$ is the 
set of edges, and $w: E \to \mathbb{R}^+$ is the cost function. The objective is to find a permutation $\sigma$ of the nodes that 
minimizes the sum:

\begin{equation}
    Cost = \sum_{i=1}^{n-1} w(v_{\sigma(i)}, v_{\sigma(i+1)}) + w(v_{\sigma(n)}, v_{\sigma(1)})
\end{equation}

\subsection{Real-life Applications}

Studying TSP has direct implications for clever logistics and transportation: even tiny improvements worth a few kilometers saved daily will 
contribute to big savings across vehicle fleets, interconnected warehouses, and regional networks. Its computational complexity, on the other 
hand, demonstrates that real applications won't use TSP as is. The preferred strategy is a combination of more adaptations of TSP: Multiple 
Salesmen Problem (mTSP), TSP with Time Windows (TSPTW), Time-dependent TSP, Dynamic TSP \cite{tsp_adaptations}. 
Other applications of TSP and its adaptations include: X-Ray crystallography, computer wiring, mask plotting in PCB production and printing 
press scheduling. \cite{matai2010traveling}
